% !TeX root = ../../main.tex
% =====================================================================
%  Herramientas de trabajo: CAD, datos y versiones
%  Responsable:
%  Última revisión:
% =====================================================================
\chapter{Herramientas de trabajo: CAD, datos y versiones}
\label{ch:herramientas_cad_datos}

\begin{objetivos}
  \item Montar el ensamblaje maestro del coche y explicar por qué ningún
        subsistema referencia directamente a otro.
  \item Definir el sistema de coordenadas del coche y la transformación al
        convenio de dinámica vehicular.
  \item Asignar correctamente un código a cualquier pieza, conjunto, componente
        comercial o utillaje del coche.
  \item Distinguir cuándo un cambio es una revisión y cuándo es una pieza nueva.
  \item Preparar un paquete de fabricación que un taller externo pueda ejecutar
        sin llamarte por teléfono.
\end{objetivos}

\fichasesion{90 minutos}{Capítulo \ref{ch:coche_como_sistema}}%
{Ordenador con el CAD del equipo instalado y acceso al servidor}

Este capítulo no va de aprender a usar el CAD. Va de las convenciones que hacen
que el trabajo de quince personas durante nueve meses siga siendo utilizable
cuando ninguna de esas quince siga en el equipo. Es, junto con el capítulo
\ref{ch:documentacion_traspaso}, el que más directamente sirve al propósito de
este libro.

\subsection*{Nuestras herramientas}

El equipo diseña en \textbf{Fusion 360}, con la intención de pasar a
\textbf{SolidWorks} más adelante. Eso condiciona tres cosas de este capítulo:

\begin{itemize}
  \item \textbf{Las convenciones de este capítulo son independientes de la
        herramienta, y eso es deliberado.} Los códigos viven en el campo
        \emph{Part Number} del componente y en el nombre del fichero, no en una
        función propia de un programa. Si un día se cambia de CAD, los códigos, la
        estructura de carpetas y el registro de piezas siguen valiendo.
  \item \textbf{Fusion resuelve por defecto parte del apartado
        \ref{sec:herramientas_cad_datos:04}}: guarda el histórico de versiones de
        cada fichero y centraliza los datos en la nube. Lo que \emph{no} resuelve
        solo es la organización: los permisos por carpeta, quién es responsable de
        qué y las congelaciones de hitos hay que montarlos igual.
  \item \textbf{La migración no es gratis.} Al pasar geometría de un CAD a otro
        por un formato neutro tipo STEP viaja el sólido, pero no el árbol de
        operaciones: en el programa de destino aparece una pieza que no se puede
        editar paramétricamente. Por eso una migración se hace al empezar un coche
        nuevo, remodelando, y no a mitad de temporada.
\end{itemize}

\pendiente{Cuando se cierre la decisión sobre SolidWorks, añadir aquí las
instrucciones concretas de cada programa: plantillas de pieza y de plano, cómo se
rellenan las propiedades y cómo se exporta la lista de materiales.}

% ---------------------------------------------------------------------
\section{Estructura del ensamblaje maestro}
\label{sec:herramientas_cad_datos:01}
\index{ensamblaje maestro}

Todo el coche cuelga de un único ensamblaje de nivel superior, y por debajo hay
exactamente un nivel de conjuntos de subsistema:

\begin{center}
\begin{minipage}{0.8\textwidth}
\ttfamily\small
FS26-GEN-A000\quad coche completo\\
\hspace*{1em}FS26-GEN-S001\quad esqueleto (referencia de todos)\\
\hspace*{1em}FS26-CHS-A000\quad conjunto chasis\\
\hspace*{1em}FS26-SUS-A000\quad conjunto suspensión\\
\hspace*{2em}FS26-SUS-A001\quad esquina delantera izquierda\\
\hspace*{2em}FS26-SUS-A002\quad esquina delantera derecha\\
\hspace*{1em}FS26-AER-A000\quad conjunto aerodinámica\\
\hspace*{1em}\ldots
\end{minipage}
\end{center}

Y una única regla, que es la más importante de todo el capítulo:

\begin{clave}
\textbf{Ningún subsistema referencia geometría de otro subsistema.} Todas las
referencias externas van contra el esqueleto, y solo contra el esqueleto. Si la
mangueta necesita saber dónde está el anclaje del trapecio, ese punto está en el
esqueleto; no se coge del chasis.
\end{clave}

El motivo es práctico: en cuanto un modelo referencia a otro, abrir una pieza
obliga a cargar medio coche, un cambio en el chasis puede romper silenciosamente
la suspensión, y dos personas no pueden trabajar a la vez. Con el esqueleto como
único punto de contacto, cada subsistema es independiente y los cambios que
afectan a todos pasan por un sitio en el que se ven.

% ---------------------------------------------------------------------
\section{Esqueleto y puntos duros}
\label{sec:herramientas_cad_datos:02}
\index{esqueleto}\index{puntos duros}

El esqueleto (\texttt{FS26-GEN-S001}) es una pieza sin volumen que contiene solo
geometría de referencia. Es el contrato entre subsistemas hecho fichero.

\subsection{Sistema de coordenadas}

Antes que nada hay que fijar el origen y la orientación de los ejes, escribirlo,
y no volver a discutirlo:

\begin{itemize}
  \item \textbf{Origen}: intersección del plano de simetría del coche, el plano
        del suelo y el eje delantero (centro de las ruedas delanteras proyectado
        al suelo), con el coche a la altura de marcha de diseño.
  \item \textbf{$X$} positivo hacia atrás. Así la batalla y las posiciones
        longitudinales son números positivos crecientes hacia la cola.
  \item \textbf{$Y$} positivo hacia la derecha del piloto.
  \item \textbf{$Z$} positivo hacia arriba.
\end{itemize}

\begin{ojo}
El convenio de ejes del CAD \textbf{no} coincide con el convenio SAE que se usa
en dinámica vehicular (véase la nomenclatura), donde $Z$ apunta hacia abajo, lo
que provoca errores de signo al pasar cargas de un sitio a otro.
La transformación entre ambos se escribe \emph{una vez}, en la nomenclatura de
este libro, y se aplica siempre desde ahí. Nunca «a ojo».
\end{ojo}

\subsection{Contenido del esqueleto}

\begin{center}
\begin{tabularx}{\textwidth}{@{}lX@{}}
\toprule
\textbf{Elemento} & \textbf{Qué incluye} \\
\midrule
Planos de referencia & Simetría, suelo, eje delantero, eje trasero \\
Cotas maestras & Batalla, vías delantera y trasera, altura de marcha de diseño \\
Centros de rueda & Los cuatro, en posición de diseño \\
Puntos duros de suspensión & Todos los anclajes, chasis y mangueta (capítulo \ref{ch:puntos_duros}) \\
Referencia del piloto & Punto del asiento, plantillas de reglamento \\
Volúmenes reservados & Espacios que un subsistema no puede invadir \\
\bottomrule
\end{tabularx}
\end{center}

Los volúmenes reservados merecen una mención: son cajas sin más significado que
«aquí no entra nadie más». Reservar desde el principio el volumen del radiador,
el del acumulador o el del recorrido del mazo evita la mitad de los conflictos
de integración del capítulo \ref{ch:coche_como_sistema}.

\subsection{Quién puede tocarlo}

El esqueleto tiene \textbf{un único responsable} con permiso de escritura ---
normalmente el responsable de integración o el de chasis. Cualquier cambio en él
afecta a todo el coche, así que se comunica al equipo entero y se anota con
fecha y motivo. Un cambio de punto duro que no se comunica se acaba descubriendo
en el montaje, con las piezas ya fabricadas.

% ---------------------------------------------------------------------
\section{Nomenclatura y numeración de piezas}
\label{sec:herramientas_cad_datos:03}
\index{nomenclatura}\index{codificación de piezas}

\subsection{Anatomía del código}

Toda pieza, conjunto, componente comprado y utillaje del coche tiene un código
con esta forma:

\begin{figure}[H]
\centering
\begin{tikzpicture}[
    font=\sffamily\footnotesize,
    bloque/.style={anchor=west,inner xsep=0pt,inner ysep=2.5pt,outer sep=0pt,
                   font=\ttfamily\LARGE},
    etiq/.style={anchor=north,align=center,text width=2.2cm,fusalPrim},
    guia/.style={fusalGris,thin}]
  % El código, encadenado para que no haya huecos entre grupos
  \node[bloque] (g1) at (0,0)      {FS26};
  \node[bloque] (g2) at (g1.east)  {-};
  \node[bloque] (g3) at (g2.east)  {SUS};
  \node[bloque] (g4) at (g3.east)  {-};
  \node[bloque] (g5) at (g4.east)  {P};
  \node[bloque] (g6) at (g5.east)  {014};

  % Etiquetas, repartidas para que no se solapen
  \node[etiq] (e1) at (-0.9,-1.5) {\textbf{Proyecto}\\coche y temporada};
  \node[etiq] (e2) at ( 1.4,-1.5) {\textbf{Subsistema}\\3 letras, lista cerrada};
  \node[etiq] (e3) at ( 3.7,-1.5) {\textbf{Tipo}\\1 letra};
  \node[etiq] (e4) at ( 6.0,-1.5) {\textbf{Secuencial}\\3 dígitos};

  \draw[guia] (g1.south) -- ++(0,-0.35) -- (e1.north);
  \draw[guia] (g3.south) -- ++(0,-0.35) -- (e2.north);
  \draw[guia] (g5.south) -- ++(0,-0.35) -- (e3.north);
  \draw[guia] (g6.south) -- ++(0,-0.35) -- (e4.north);
\end{tikzpicture}
\caption{Anatomía del código de pieza. Trece caracteres, legibles por teléfono y
ordenables en el explorador de archivos.}
\label{fig:anatomia_codigo}
\end{figure}

\begin{clave}
El código \textbf{identifica}, no describe. Nunca se codifica en él nada que
pueda cambiar sin que cambie la pieza: material, proveedor, masa, autor o fecha.
Todo eso vive en las propiedades del modelo, donde se puede corregir sin
renumerar medio coche.
\end{clave}

\subsection{Códigos de subsistema}

Lista cerrada. Si hace falta uno nuevo, se añade aquí y se comunica; no se
improvisa.

\begin{center}
\begin{tabularx}{\textwidth}{@{}llX@{}}
\toprule
\textbf{Código} & \textbf{Subsistema} & \textbf{Alcance} \\
\midrule
\ttfamily GEN & General        & Esqueleto, conjunto general, tornillería propia, misceláneo \\
\ttfamily CHS & Chasis         & Estructura primaria, anclajes, atenuador, mamparos \\
\ttfamily SUS & Suspensión     & Trapecios, manguetas, bujes, balancines, tirantes, muelles, amortiguadores, barras \\
\ttfamily DIR & Dirección      & Cremallera, columna, volante, bieletas \\
\ttfamily FRE & Frenos         & Discos, pinzas, bombas, pedalera de freno, conductos \\
\ttfamily RUE & Ruedas         & Llantas, neumáticos, tuercas y centros de rueda \\
\ttfamily AER & Aerodinámica   & Alerones, fondo, difusor, carrocería, soportes \\
\ttfamily REF & Refrigeración  & Radiadores, ventiladores, conductos, bomba, circuito \\
\ttfamily TRA & Transmisión    & Reductora, diferencial, palieres, cadena o correa, resguardos \\
\ttfamily MOT & Motor          & Grupo motopropulsor y sus soportes \\
\ttfamily ACU & Acumulador     & Celdas, contenedor, refrigeración de batería (vehículos eléctricos) \\
\ttfamily ELE & Electrónica    & Placas, cajas, sensores, mazos \\
\ttfamily ERG & Ergonomía      & Asiento, pedalera, reposacabezas, arneses, cockpit \\
\bottomrule
\end{tabularx}
\end{center}

\pendiente{Revisar esta lista con los responsables de subsistema y cerrarla antes
de empezar a modelar. Cambiarla a mitad de temporada cuesta mucho más de lo que parece.}

\subsection{Letras de tipo}

\begin{center}
\begin{tabularx}{\textwidth}{@{}clX@{}}
\toprule
\textbf{Letra} & \textbf{Tipo} & \textbf{Uso} \\
\midrule
\ttfamily A & Conjunto & Ensamblajes. El \ttfamily A000 \normalfont de cada subsistema es su conjunto principal \\
\ttfamily P & Pieza    & Piezas que fabricamos nosotros o encargamos fabricar según plano \\
\ttfamily C & Comercial & Componentes comprados de catálogo \\
\ttfamily U & Utillaje & Moldes, plantillas, útiles de soldadura, bancos de ensayo \\
\ttfamily S & Simulación & Geometría de referencia, maquetas y volúmenes reservados. No se fabrica \\
\bottomrule
\end{tabularx}
\end{center}

El secuencial son tres dígitos, empieza en \texttt{001} y es independiente para
cada combinación de subsistema y tipo. Es decir, pueden coexistir
\texttt{FS26-SUS-P014} y \texttt{FS26-SUS-C014}: son cosas distintas y no se
confunden porque la letra los separa.

\subsection{Las seis reglas de numeración}

\begin{enumerate}
  \item \textbf{Un número no se reutiliza nunca}, ni aunque la pieza se abandone.
        Reutilizarlo hace que dos documentos distintos digan cosas distintas del
        mismo código.
  \item \textbf{Izquierda y derecha son dos códigos distintos.} Aunque sean
        espejo, son dos piezas que se fabrican, se compran y se cuestan por
        separado, y así deben aparecer en el BOM. La relación de simetría se anota
        en las propiedades.
  \item \textbf{Los componentes comerciales también llevan código nuestro}, del
        tipo \texttt{C}. La referencia del fabricante va en las propiedades: así,
        cambiar de proveedor no obliga a tocar el diseño.
  \item \textbf{Si la pieza nueva es intercambiable con la vieja, es una
        revisión.} Si no lo es, es un código nuevo. Esta es la única prueba que
        hay que aplicar, y se aplica preguntando: ¿puedo montar la nueva donde
        estaba la vieja y que el coche funcione igual?
  \item \textbf{El código se asigna al crear el fichero}, no al terminarlo. Una
        pieza sin código no existe y no se puede referenciar.
  \item \textbf{El registro de códigos es único y compartido.} Una hoja con una
        fila por código, y quien asigna el siguiente lo anota en el momento. Sin
        ese registro aparecen números duplicados en cuanto trabajan dos personas
        a la vez.
\end{enumerate}

\subsection{Nombres de fichero}

El nombre del fichero es el código, un guion bajo, y una descripción corta:

\begin{center}
\ttfamily\large FS26-SUS-P014\_trapecio-del-sup-izq.\textrm{\itshape ext}
\end{center}

Reglas de la parte descriptiva:

\begin{itemize}
  \item Minúsculas, sin acentos, sin eñes, sin espacios. Separador: guion.
  \item Corta: unos 30 caracteres. Es una etiqueta, no una frase.
  \item Abreviaturas fijas para la posición: \texttt{del}/\texttt{tra},
        \texttt{sup}/\texttt{inf}, \texttt{izq}/\texttt{der}, en ese orden.
\end{itemize}

\begin{ojo}
Lo que \textbf{nunca} aparece en el nombre de un fichero:

\begin{center}
\ttfamily\small
trapecio\_v2.ext \quad trapecio\_FINAL.ext \quad trapecio\_final\_bueno.ext\\
trapecio\_2026-03-14.ext \quad trapecio\_JL.ext \quad trapecio\_OK\_este\_si.ext
\end{center}

\noindent
La versión, la fecha y el autor los guarda el sistema de archivos y el control de
versiones, y los guarda bien. Escritos en el nombre, lo único que consiguen es que
dentro de unos meses nadie sepa cuál es el fichero bueno.
\end{ojo}

\subsection{Revisiones y estados}
\index{revisiones}

La revisión es una letra: \texttt{A} para la primera versión liberada,
\texttt{B} para la siguiente, y así. \textbf{Solo hay revisiones después de la
primera liberación}: mientras la pieza se está diseñando no se revisa, se
modifica.

Cada pieza tiene además un estado, y el estado decide qué se puede hacer con ella:

\begin{center}
\begin{tabularx}{\textwidth}{@{}lXl@{}}
\toprule
\textbf{Estado} & \textbf{Significado} & \textbf{¿Se fabrica?} \\
\midrule
En diseño   & Trabajo en curso. Puede cambiar en cualquier momento & No \\
En revisión & Terminada, pendiente de que la revise otra persona & No \\
Liberada    & Revisada y aprobada. Plano firmado & \textbf{Sí} \\
Obsoleta    & Sustituida o abandonada. Se conserva, no se borra & No \\
\bottomrule
\end{tabularx}
\end{center}

\begin{clave}
Al taller solo va material \textbf{liberado}, y liberar exige que lo haya revisado
alguien que no sea el autor. Cuesta unos minutos y detecta errores que el autor ya
no ve.
\end{clave}

Cuando una pieza liberada cambia, la revisión sube una letra, se anota en la
tabla de revisiones del plano (qué cambió y por qué) y \emph{se avisa a quien
tenga la revisión anterior en la mano}: si no, seguirá trabajando con un plano
obsoleto sin saberlo.

\subsection{Propiedades obligatorias del modelo}
\index{propiedades del modelo}

Cada pieza lleva rellenas estas propiedades personalizadas. No es burocracia: son
las que alimentan automáticamente el cajetín del plano, la lista de materiales y
el informe de costes del capítulo \ref{ch:evento_costes}. Rellenarlas al crear la
pieza cuesta un minuto; reconstruirlas en junio para el informe de costes cuesta
semanas.

\begin{center}
\begin{tabularx}{\textwidth}{@{}lXl@{}}
\toprule
\textbf{Propiedad} & \textbf{Contenido} & \textbf{Obligatoria} \\
\midrule
\ttfamily CODIGO        & Código completo de la pieza. En Fusion, el campo
                          \emph{Part Number} del componente & Sí \\
\ttfamily DENOMINACION  & Nombre en castellano, como aparece en el plano.
                          En Fusion, \emph{Description} & Sí \\
\ttfamily ESTADO        & En diseño / en revisión / liberada / obsoleta & Sí \\
\ttfamily REVISION      & Letra de revisión & Sí \\
\ttfamily AUTOR         & Quien la diseña & Sí \\
\ttfamily REVISOR       & Quien la aprueba al liberarla & Al liberar \\
\ttfamily FECHA         & Fecha de liberación & Al liberar \\
\ttfamily MATERIAL      & Designación normalizada, no «aluminio» & Sí \\
\ttfamily MASA          & Calculada por el CAD & Automática \\
\ttfamily PROCESO       & Mecanizado, soldadura, laminado, impresión\ldots & Sí \\
\ttfamily CANTIDAD      & Unidades por coche & Sí \\
\ttfamily PROVEEDOR     & Solo para comerciales y subcontratados & Si aplica \\
\ttfamily REF\_PROVEEDOR & Referencia de catálogo del fabricante & Si aplica \\
\ttfamily TRATAMIENTO   & Térmico o superficial, si lo lleva & Si aplica \\
\ttfamily SUBSISTEMA    & Para agrupar el BOM & Sí \\
\bottomrule
\end{tabularx}
\end{center}

\subsection{Carpetas}

\begin{center}
\begin{minipage}{0.8\textwidth}
\ttfamily\small
FS26\_CAD/\\
\hspace*{1em}00\_GEN/\quad esqueleto, conjunto general, biblioteca de tornillería\\
\hspace*{1em}01\_CHS/\\
\hspace*{2em}Piezas/\\
\hspace*{2em}Conjuntos/\\
\hspace*{2em}Comerciales/\\
\hspace*{2em}Planos/\\
\hspace*{2em}Fabricacion/\quad paquetes listos para taller (PDF + STEP + DXF)\\
\hspace*{1em}02\_SUS/\quad\textrm{(misma estructura)}\\
\hspace*{1em}\ldots\\
\hspace*{1em}98\_Obsoleto/\quad nada se borra, todo se mueve aquí\\
\hspace*{1em}99\_Congelaciones/\quad copias de solo lectura de cada hito
\end{minipage}
\end{center}

% ---------------------------------------------------------------------
\section{Control de versiones y copias de seguridad}
\label{sec:herramientas_cad_datos:04}
\index{copias de seguridad}

El CAD de un coche es el activo más valioso que produce el equipo y el más fácil
de perder. Las reglas son pocas y no admiten excepción:

\begin{enumerate}
  \item \textbf{Hay un único original, y está en el servidor del equipo.} Lo que
        hay en un portátil es una copia de trabajo, nunca la referencia. Un
        proyecto que vive en el ordenador de alguien se pierde el día que esa
        persona se va, y se va siempre.
  \item \textbf{Una carpeta, un responsable con permiso de escritura.} Los demás
        leen. Si alguien necesita modificar una pieza que no es suya, se la pide
        a su responsable.
  \item \textbf{Nadie abre el mismo fichero desde dos sitios.} Si el CAD no tiene
        gestor de datos que bloquee ficheros, se suple con una hoja compartida de
        control: quién tiene qué abierto y desde cuándo. Es rudimentario y funciona.
  \item \textbf{Copia de seguridad automática diaria}, conservando al menos
        treinta días. Automática: lo que depende de que alguien se acuerde no es
        una copia de seguridad.
  \item \textbf{Congelación en cada hito.} Al cerrar el diseño de un subsistema se
        guarda una copia completa de \emph{solo lectura} en
        \texttt{99\_Congelaciones/}, con el nombre
        \texttt{FS26\_congelacion\_2026-02-15}. Es la única forma de poder
        contestar dentro de un año a «¿qué mandamos a fabricar exactamente?».
  \item \textbf{Al terminar la temporada, archivo completo del coche tal y como
        se fabricó}, incluyendo planos, paquetes de fabricación y el registro de
        códigos. Ese archivo es el que se consultará dentro de tres años, y por
        entonces no quedará nadie a quien preguntar.
\end{enumerate}

\begin{ojo}
\begin{itemize}
  \item Trabajar con el CAD en una carpeta sincronizada a la nube y abierta desde
        dos ordenadores a la vez. Las herramientas de sincronización resuelven los
        conflictos duplicando ficheros, y en un ensamblaje eso rompe las
        referencias sin avisar.
  \item Confundir sincronización con copia de seguridad. Si borras una carpeta,
        la nube sincroniza el borrado obedientemente.
  \item Guardar las congelaciones con permiso de escritura. Acaban modificadas.
\end{itemize}
\end{ojo}

% ---------------------------------------------------------------------
\section{Planos, acotación y entrega a taller}
\label{sec:herramientas_cad_datos:05}
\index{planos}

\subsection{El plano es el contrato}

Lo que no está en el plano no existe. El taller no sabe lo que tenías en la
cabeza, no va a abrir tu modelo y, si algo es ambiguo, elegirá la interpretación
más barata de fabricar --- que es lo correcto por su parte.

\begin{itemize}
  \item \textbf{Formato A3} por defecto, con cajetín rellenado automáticamente
        desde las propiedades del modelo.
  \item \textbf{Acotación funcional}: se acota desde la superficie que define la
        función de la pieza, no desde la esquina que resultaba cómoda. Las
        referencias de acotación deben coincidir con las de montaje.
  \item \textbf{Tolerancia general} indicada en el cajetín (por ejemplo
        ISO 2768-mK) y tolerancias estrechas \emph{solo} donde importan. Cada
        tolerancia apretada es dinero y tiempo de máquina; si están todas
        apretadas, en realidad no has decidido ninguna.
  \item \textbf{Acabado superficial} solo donde tenga función.
  \item \textbf{Tabla de revisiones} en el plano, con fecha y motivo de cada una.
\end{itemize}

\subsection{Paquete de fabricación}

Cada pieza que va a taller se entrega como un paquete completo en
\texttt{Fabricacion/}, con el código y la revisión en el nombre:

\begin{center}
\begin{tabularx}{\textwidth}{@{}lX@{}}
\toprule
\textbf{Fichero} & \textbf{Para qué} \\
\midrule
\ttfamily \ldots\_planoB.pdf  & El plano firmado. Es el documento contractual \\
\ttfamily \ldots\_revB.step   & Geometría neutra para programar la máquina \\
\ttfamily \ldots\_revB.dxf    & Solo para corte plano: láser, agua, plegado \\
\bottomrule
\end{tabularx}
\end{center}

\begin{clave}
A un proveedor externo se le envía \textbf{siempre} STEP y PDF, nunca el fichero
nativo del CAD. El nativo lleva el histórico de operaciones, las relaciones con el
resto del coche y a menudo información que no queremos repartir; además obliga al
proveedor a tener nuestro mismo programa y versión.
\end{clave}

% ---------------------------------------------------------------------
%  Cierre del capítulo
% ---------------------------------------------------------------------

\begin{caso}
\redactar{Rellenar con la situación real: qué CAD usamos, dónde vive el modelo
hoy, si existe esqueleto, qué convenio de códigos veníamos usando y qué se ha
perdido de temporadas anteriores. Conviene ser explícito sobre lo que no
tenemos: es lo que justifica adoptar este capítulo.}
\end{caso}

\begin{ojo}
\begin{itemize}
  \item Empezar a modelar antes de tener el esqueleto. Se acaba rehaciendo todo.
  \item Referenciar geometría de otro subsistema «solo esta vez».
  \item Dejar las propiedades del modelo para el final: en junio, cuando hay que
        entregar el informe de costes, ya no hay quien las reconstruya.
  \item Numerar por orden de fabricación en lugar de por subsistema y tipo.
  \item Enviar a taller un plano de una pieza que no está liberada.
  \item Tolerar en el plano cotas que no se pueden medir con lo que hay en el
        equipo. Si no la puedes verificar, no la puedes exigir.
\end{itemize}
\end{ojo}

\begin{entregable}
La \textbf{guía de estándares CAD del equipo}, que es este capítulo con las
decisiones concretas tomadas:

\begin{enumerate}
  \item Tabla de códigos de subsistema cerrada y aprobada por los responsables.
  \item Fichero de esqueleto creado, con el sistema de coordenadas documentado y
        la transformación al convenio SAE escrita.
  \item Plantilla de pieza y de plano con las propiedades obligatorias y el
        cajetín ya enlazado.
  \item Estructura de carpetas creada en el servidor, con permisos asignados.
  \item Registro de códigos abierto, con un responsable de asignación.
\end{enumerate}
\end{entregable}

\begin{juez}
\begin{enumerate}
  \item ¿Cómo os aseguráis de que el modelo que está en el taller es el mismo que
        el que habéis analizado?
  \item Enséñame el plano de esta pieza. ¿Por qué está acotada desde ahí?
  \item ¿Por qué esta tolerancia es tan estrecha? ¿Cómo la verificáis?
  \item ¿Qué pasa cuando cambiáis un punto duro de la suspensión? ¿Quién se entera?
  \item Si tu compañero de subsistema deja el equipo mañana, ¿podrías seguir su
        trabajo?
\end{enumerate}
\end{juez}
